\documentclass[aspectratio=169, ignorenonframetext]{beamer}
\usepackage{pgf}

\mode<article>{
  \usepackage{fullpage}

}

\mode<presentation>
{
%   \usetheme{Madrid}
%   \useinnertheme{circles}
  \setbeamertemplate{navigation symbols}{}
%   \usetheme[hideothersubsections, width=2.4cm]{Hannover}
%   \usetheme{Antibes}
%   \usetheme{Montpellier}
%   \usetheme{Singapore}
  \usetheme{Darmstadt}
%   \usecolortheme{seahorse}
%   \usecolortheme{orchid}
  \setbeamertemplate{footline}
  {%
%     \begin{beamercolorbox}{section in head/foot}
%       \usebeamercolor{bg}
       \hskip 1em \footnotesize \insertframenumber{} / \inserttotalframenumber%\hskip 41em \includegraphics[height = .5cm]{../../../../bilder/cc_by-nc_eu.png}
% % cc_by-nc_eu.png: 403x141 px, 72dpi, 14.22x4.97 cm, bb=0 0 403 141
%
% % bwslogo_3.png: 476x392 px, 300dpi, 4.03x3.32 cm, bb=0 0 114 94
%
%       %\hskip 5em
%       %       \input{../bilder/cc_by.png}
%       %\includegraphics[height=.5cm]{../../../../bilder/bwslogo_3.png}
%     \end{beamercolorbox}%
  }
  \usepackage{beamerfoils}
}
\usepackage[ngerman]{babel}
\usepackage[utf8]{inputenc}
\usepackage{times}
\usepackage[T1]{fontenc}
\usepackage{pgf}
\usepackage{pgfpages}
\usepackage{eurosym}
\usepackage{graphicx}
\usepackage{amsmath}
\usepackage[siunitx,european]{circuitikz}
\usepackage{ulem}
\usepackage{listings}
\usepackage{svg}
%
\lstset{numbers=left, numberstyle=\tiny, stepnumber=2, numbersep=5pt, language = C++, alsolanguage=XML}
% \MyLogo{\includegraphics[height=1cm]{../../../../bilder/bwslogo_3.png}}
% % \includegraphics{../../bilder/bwslogo_3.png}
% % bwslogo_3.png: 476x392 px, 300dpi, 4.03x3.32 cm, bb=
%

\only<presentation>{
  \usepackage{hyperref}
}


\title{Maschenstromverfahren  Aufgabe 1}

\institute[BWS Hofheim]{Brühlwiesenschule, Hofheim}
\author{Thomas Maul}

\titlegraphic{Für eigene Teile gilt: \includegraphics[height=1cm]{../../../../../../bilder/cc_by-nc-sa.eu.png}}

\begin{document}
  \begin{frame}<beamer>
    \titlepage
    % \hyperlink{Teil_2}{\beamerbutton{Go part 2}}
  \end{frame}



\begin{frame}{LGS aufstellen (Schritt 1)}
$\left( \begin{array}{lll}
\text{R}_1+\text{R}_2+\text{R}_3 & & \\
& \text{R}_3+\text{R}_4+\text{R}_5 & \\
&&\text{R}_2 + \text{R}_4 + \text{R}_6 \\
\end{array}\right)
% *
% \left(\begin{array}{c}
%   \text{I}_{M1}\\
%   \text{I}_{M2}\\
%   \text{I}_{M3}\\
%      \end{array} \right)
%      =
% \left(\begin{array}{c}
%   \text{U}_{01}-\text{U}_{02}\\
%   -\text{U}_{02}+\text{U}_{03}\\
%   0\\
%      \end{array} \right)
$
\end{frame}

\begin{frame}{LGS aufstellen (Schritt 2)}
$\left( \begin{array}{lll}
\text{R}_1+\text{R}_2+\text{R}_3 & & \\
& \text{R}_3+\text{R}_4+\text{R}_5 & \\
&&\text{R}_2 + \text{R}_4 + \text{R}_6 \\
\end{array}\right) *
\left(\begin{array}{c}
  \text{I}_{M1}\\
  \text{I}_{M2}\\
  \text{I}_{M3}\\
     \end{array} \right)
%      =
% \left(\begin{array}{c}
%   \text{U}_{01}-\text{U}_{02}\\
%   -\text{U}_{02}+\text{U}_{03}\\
%   0\\
%      \end{array} \right)
$
\end{frame}

\begin{frame}{LGS aufstellen (Schritt 3)}
$\left( \begin{array}{lll}
\text{R}_1+\text{R}_2+\text{R}_3 & & \\
& \text{R}_3+\text{R}_4+\text{R}_5 & \\
&&\text{R}_2 + \text{R}_4 + \text{R}_6 \\
\end{array}\right) *
\left(\begin{array}{c}
  \text{I}_{M1}\\
  \text{I}_{M2}\\
  \text{I}_{M3}\\
     \end{array} \right)
     =
\left(\begin{array}{c}
  \text{U}_{01}-\text{U}_{02}\\
  -\text{U}_{02}+\text{U}_{03}\\
  0\\
     \end{array} \right)
$
\end{frame}

\begin{frame}{LGS aufstellen (Schritt 4)}
$\left( \begin{array}{ccc}
\text{R}_1+\text{R}_2+\text{R}_3 &\text{R}_3 & \\
\text{R}_3& \text{R}_3+\text{R}_4+\text{R}_5 & \\
&&\text{R}_2 + \text{R}_4 + \text{R}_6 \\
\end{array}\right) *
\left(\begin{array}{c}
  \text{I}_{M1}\\
  \text{I}_{M2}\\
  \text{I}_{M3}\\
     \end{array} \right)
     =
\left(\begin{array}{c}
  \text{U}_{01}-\text{U}_{02}\\
  -\text{U}_{02}+\text{U}_{03}\\
  0\\
     \end{array} \right)
$
\end{frame}

\begin{frame}{LGS aufstellen (Schritt 5)}
$\left( \begin{array}{ccc}
\text{R}_1+\text{R}_2+\text{R}_3 &\text{R}_3 & \text{R}_2\\
\text{R}_3& \text{R}_3+\text{R}_4+\text{R}_5 & -\text{R}_4\\
\text{R}_2&-\text{R}_4&\text{R}_2 + \text{R}_4 + \text{R}_6 \\
\end{array}\right) *
\left(\begin{array}{c}
  \text{I}_{M1}\\
  \text{I}_{M2}\\
  \text{I}_{M3}\\
     \end{array} \right)
     =
\left(\begin{array}{c}
  \text{U}_{01}-\text{U}_{02}\\
  -\text{U}_{02}+\text{U}_{03}\\
  0\\
     \end{array} \right)
$
\end{frame}

\begin{frame}{Werte einsetzen (I)}
$\left( \begin{array}{ccc}
330\Omega+470\Omega + 270\Omega & 270\Omega& 470\Omega\\
270\Omega & 270\Omega+680\Omega+330\Omega & -680\Omega\\
470\Omega &-680\Omega&470\Omega + 680\Omega + 470\Omega \\
\end{array}\right) *
\left(\begin{array}{c}
  \text{I}_{M1}\\
  \text{I}_{M2}\\
  \text{I}_{M3}\\
     \end{array} \right)
     =
\left(\begin{array}{c}
  24\,V-18\,V\\
  -18\,V+16\,V\\
  0\\
     \end{array} \right)
$
\end{frame}

\begin{frame}{Werte einsetzen (II)}
$\left( \begin{array}{ccc}
1070\,\Omega & 270\,\Omega& 470\,\Omega\\
270\,\Omega & 1280\,\Omega & -680\,\Omega\\
470\,\Omega &-680\,\Omega&1620\,\Omega \\
\end{array}\right) *
\left(\begin{array}{c}
  \text{I}_{M1}\\
  \text{I}_{M2}\\
  \text{I}_{M3}\\
     \end{array} \right)
     =
\left(\begin{array}{c}
  6V\\
  -2V\\
  0\\
     \end{array} \right)
$
\end{frame}

\begin{frame}{Eingabe in Octave}
 R = [1070, 270, 470;\\
     270, 1280, 680;\\
     470, 680, 1620];\\[\baselineskip]
u = [6;\\
     2;\\
     0];\\[\baselineskip]
I = R \textbackslash{} u;\\[\baselineskip]
disp('Ströme in mA:');
disp(I * 1000);
\end{frame}

\begin{frame}{Ströme}
  Maschenströme\\
 $\begin{array}{llcr}
  \text{I}_{M1} &= I_1 &=& 6,31\ mA\\
  \text{I}_{M2} &= I_5 &=& 1,55\ mA\\
  \text{I}_{M3} &= I_6 &=& -2,48\ mA\\
 \end{array}$
% \end{frame}
\pause
% \begin{frame}{
 \\[\baselineskip]Übrige Ströme\\
 $\begin{array}{llcrr}
   \text{I}_{2} &= \text{I}_{M1} + \text{I}_{M3} & = & 6,31\ mA 2,48\ mA & = 3,83\ mA\\
  \text{I}_{3} &= \text{I}_{M1} + \text{I}_{M2} &=& 6,31\ mA + 1,55\ mA &= 7,86\ mA\\
  \text{I}_{4} &=\text{I}_{M2} - \text{I}_{M3} &=& 1,55\ mA + 2,48\ mA &= 4,03\ mA
 \end{array}$
\end{frame}

\end{document}
